\documentclass[11pt]{article}
\usepackage[a4paper,margin=25mm]{geometry}
\usepackage{amsmath,amssymb,booktabs,longtable,hyperref}
\hypersetup{hidelinks}
\renewcommand{\theequation}{\arabic{equation}c}
\newcommand{\Lc}{\textbf{(c)}}
\newcommand{\Ld}{\textbf{(d)}}
\title{Projected collective measure: exact Wronskians,\\tail cancellation and a boundary-domain obstruction}
\author{Independent project audit}
\date{4 November 2026}
\begin{document}
\maketitle

\paragraph{Concrete result.}
The nonzero potential-tail cross terms do \emph{not} change the finite-dimensional
collective symplectic volume. Their contribution cancels from the collective Wronskian determinant, as
verified directly on the exact two-soliton background during the collision
\Lc. There is nevertheless a precise obstruction to treating this result as a
projected field determinant: the natural symplectic projector onto the full
soliton family can send a localized fluctuation to a variation with nonzero
potential endpoint. It therefore does not preserve a fluctuation domain with
fixed potential endpoints \Lc. A boundary-domain prescription and a
finite-time determinant construction remain necessary \Ld.

This continues \texttt{audit\_determinant\_foundations.tex}; it neither selects
a field regularisation nor assigns an exchange phase. The repository good
equation, canonical form $\Omega=\int d\pi\wedge d\phi$ and potential gauge
$\phi(-\infty)=0$ are retained. The calculations are for solution-parameter
variations in that gauge. Every numbered formula carries a claim-(c) suffix and is checked in
\texttt{scripts/audit\_\allowbreak projected\_\allowbreak measure.py}, with exact SymPy algebra and
55-digit mpmath quadrature. Numerical results are saved in
\texttt{scripts/audit\_\allowbreak projected\_\allowbreak measure\_\allowbreak results.json}.

\section{The collective Wronskian has a factorised determinant}
Let $a_j=\alpha_j$, $d_j=dP_j/da_j$, and order the tangent coordinates as
$z=(a_1,a_2,X_1,X_2)$. For separated solitons with particle~1 on the left,
\begin{equation}\label{eq:W}
 W_{AB}=\Omega(\partial_{z_A},\partial_{z_B})
 =\begin{pmatrix}C&D\\-D&0\end{pmatrix},\qquad
 D=\operatorname{diag}(d_1,d_2),\quad
 C=\begin{pmatrix}0&c\\-c&0\end{pmatrix},\quad
 c=-M'_1\Pi'_2.
\end{equation}
Here $M=2\sin a$, $\Pi=-2\epsilon\sin a\cos a$ and
$P=(2/3)\epsilon\sin^3a\cos a$. The earlier audit directly checked the
mixed entry by spatial integration. We now retain it in the full matrix.
For arbitrary $c$, not only its Boussinesq value, exact matrix algebra gives
\begin{equation}\label{eq:det}
 \det W=(d_1d_2)^2,\qquad
 W^{-1}=\begin{pmatrix}0&-D^{-1}\\D^{-1}&D^{-1}CD^{-1}\end{pmatrix}.
\end{equation}
Hence the positive Liouville density is
\begin{equation}\label{eq:volume}
 d\mu_{\rm coll}=|d_1d_2|\,da_1\,da_2\,dX_1\,dX_2.
\end{equation}
There is no tail-dependent factor in this density. This resolves a specific
concern raised in the earlier audit: individual mixed Wronskians need not
vanish for their determinant to factorise.

The established Darboux coordinates make the cancellation explicit:
\begin{equation}
 X'_1=X_1,\qquad X'_2=X_2+\gamma_2 M_1,\qquad
 \gamma_2=\Pi'_2/P'_2.
\end{equation}
Their tangent Jacobian is triangular with unit determinant. Pulling
$\sum_jdP_j\wedge dX'_j$ back gives \eqref{eq:W}; derivatives of $\gamma_2$
contribute only to a wedge with itself. The symbolic test verifies the full
matrix pullback, not merely its determinant. Switching the spatial ordering
changes $c$ and the corresponding shear, but leaves \eqref{eq:det} and
\eqref{eq:volume} unchanged. These formulas require $d_1d_2\ne0$ and do not
extend through the cusp by division.

\section{Exact collision-time verification}
Use the actual tau solution
\begin{equation}
 \tau=1+E_1+E_2+AE_1E_2,\quad
 E_j=e^{k_j(x-v_jt-x_{0j})},\quad
 \phi=2\partial_x\log\tau,\quad \pi=\partial_t\phi,
\end{equation}
with the repository's real KP cross-ratio $A>0$. The script differentiates
$\phi$ and $\pi$ with respect to $(a_1,a_2,x_{01},x_{02})$ analytically,
then integrates every entry of their $4\times4$ Wronskian matrix. The
Hirota parameters $x_{0j}$ are not confused with both sets of physical
asymptotic centers.

For $v_1>v_2$, the incoming centers are
$X_1^{\rm in}=x_{01}$ and $X_2^{\rm in}=x_{02}-\log A/k_2$; the outgoing
ones are $X_1^{\rm out}=x_{01}-\log A/k_1$ and
$X_2^{\rm out}=x_{02}$. Thus, in the fixed Hirota parameter basis, the
mixed entry becomes
\begin{align}\label{eq:c_tau}
 c_\tau
 &=\frac{d_2}{k_2}\partial_{a_1}\log A-M'_1\Pi'_2\nonumber\\
 &=-\frac{d_1}{k_1}\partial_{a_2}\log A+\Pi'_1M'_2.
\end{align}
The other blocks in \eqref{eq:W} remain $D$ and zero. The equality of the
two expressions is checked numerically along with the full spatial integral.
At $(a_1,a_2,\epsilon_1,\epsilon_2)=(0.3,0.7,+1,+1)$, all entries were
checked at $t=-8,0,8$; the maximum discrepancy is below $1.3\times10^{-55}$.
A head-on point $(0.25,0.45,+1,-1)$ at $t=0$ gives an error below
$4.1\times10^{-56}$. In every case the determinant divided by
$(d_1d_2)^2$ is one to the stated precision \Lc. These are observed
residuals; the assertions require entry errors below $10^{-43}$ and a
determinant-ratio error below $10^{-42}$.

\paragraph{Why this is conserved despite the tails.}
For potential-field linearised solutions $f,g$, write
$L=\partial_x^4-\partial_x^2+6\partial_xu\partial_x$.
A direct symbolic identity gives
\begin{align}
 \partial_x B(f,g)&=fLg-gLf,\\
 B(f,g)&=fg'''-f'g''+f''g'-f'''g
                  +(6u-1)(fg'-gf').
\end{align}
Since $f_{tt}=-Lf$, the time derivative of
$\int(f_tg-g_tf)dx$ is the endpoint difference of $B$.
Soliton-family potential variations are bounded at spatial infinity and
all their spatial derivatives decay there; therefore this boundary flux
vanishes. The integrated Wronskian is conserved even though the amplitude
variations themselves have constant tails. This supplies the missing direct
justification for comparing the finite-time matrix to its asymptotic value.
It proves a property of the finite collective solution space, not
completeness of the continuum field modes.

\section{An explicit symplectic projector and its limitation}
Let $Z_A$ denote the collective tangent solutions at a fixed time, and suppose
a variation $\eta$ belongs to a domain on which all pairings converge. Define
\begin{equation}\label{eq:projector}
 b_A=\Omega(Z_A,\eta),\qquad
 \mathcal P\eta=Z_A(W^{-1})^{AB}b_B,\qquad
 \mathcal Q=1-\mathcal P.
\end{equation}
Matrix multiplication gives
\begin{equation}
 \mathcal P^2=\mathcal P,\quad \mathcal PZ_A=Z_A,\quad
 \Omega(Z_A,\mathcal Q\eta)=0.
\end{equation}
Replacing $Z$ by $ZR$ for an invertible collective-coordinate change leaves
$\mathcal P$ unchanged, since $W$ changes to $R^TWR$. The script verifies
these identities in a six-dimensional canonical embedding with arbitrary
$d_1,d_2,c$, including the Darboux shear. Thus tail cross terms can be
incorporated in a concrete projector, rather than discarded.

\paragraph{A localized counterexample to domain preservation.}
Already for one soliton, take the localized initial-data perturbation
$\eta=(\eta_\phi,\eta_\pi)=(0,u)$, where
$u=(k^2/2)\operatorname{sech}^2[k(x-X)/2]$.
Since $\phi_X=-u$ and $\phi_a(+\infty)=2\cos a$, one finds
\begin{align}
 b_X&=\int u^2dx=\frac23 k^3,\\
 (\mathcal P\eta)^a&=-\frac{2k^3}{3P'},\qquad
 (\mathcal P\eta)_\phi(+\infty)
 =-\frac{4\cos a\,k^3}{3P'}\ne0.
\end{align}
The integral and projection are checked exactly. The original perturbation
has zero potential variation at both endpoints; its projected part and its
complement have compensating nonzero right tails. Consequently neither is
in the same fixed-endpoint domain. This is an explicit obstruction, not
merely the observation that $\phi_a$ is not square-integrable.

An enlarged space allowing such endpoint variations can accommodate
\eqref{eq:projector}. A fixed-total-$M$ space instead restricts the collective
family. On that constraint,
\begin{equation}
 da_2=-\frac{M'_1}{M'_2}da_1,\qquad
 \ker\Omega\big|_{dM_{\rm tot}=0}
 =\operatorname{span}\left\{
 \frac{M'_1}{P'_1}\partial_{X_1}
 +\frac{M'_2}{P'_2}\partial_{X_2}\right\}.
\end{equation}
The three-dimensional pullback has rank two, as checked symbolically.
It therefore cannot simply inherit the four-dimensional measure
\eqref{eq:volume}. Whether the quantum problem includes boundary variables,
uses a constrained reduction, or has an equivalent scattering-state
construction is not fixed by this calculation \Ld. No such choice is made
here.

\section{What the measure result does not determine}
\paragraph{Solution variations are not automatically determinant zero modes.}
The tangent fields solve the linearised equation, but they do not generally
satisfy the time-endpoint conditions used to define a propagator determinant.
A finite-dimensional benchmark makes the distinction precise. The free
collective equation has solutions $1$ and $t$, but neither is a nontrivial
solution with zero values at both finite endpoints. For the elementary
Dirichlet matrix $K_N=\operatorname{tridiag}(-1,2,-1)$,
\begin{equation}
 \det\left(\frac{m}{h}K_N\right)
 =(N+1)\left(\frac{m}{h}\right)^N\ne0.
\end{equation}
Laplace expansion gives $D_N=2D_{N-1}-D_{N-2}$ with $D_0=1,D_1=2$;
substituting $D_N=N+1$ proves the formula by induction. The recurrence and
base cases are checked symbolically, alongside direct determinants for
$N=1,\ldots,6$. This is a finite matrix diagnostic, not a
proposed field regularisation. ``Remove the zero modes'' must specify which
operator domain and which boundary problem are meant. A conserved solution
Wronskian is not a proof that the finite-time operator has a kernel.

\paragraph{Volume and endpoint maps do not fix the propagation phase.}
A second exact control is the unit-frequency harmonic oscillator. Its
classical symplectic flow returns to the identity at time $2\pi$, preserves
volume throughout, and has the same endpoint Jacobian as no evolution.
But the ground-state wavefunction $e^{-x^2/2}$ satisfies
$(-\partial_x^2+x^2)\psi/2=\psi/2$, so its phase after that evolution is
$e^{-i\pi}=-1$ \Lc. The script verifies both identities. Thus identical
symplectic volumes and endpoint maps do not specify the phase carried by
the evolution. This does not assign a Boussinesq Maslov index; it shows why
\eqref{eq:det} cannot do so.

\paragraph{The remaining determinant task.}
The exact result now established is factorisation of the collective
\emph{symplectic} volume and of the determinant of its solution Wronskian.
A Faddeev--Popov measure, a constrained Gaussian functional integral and a
spectral determinant depend on more data. One must give a common spatial
and temporal domain, identify permissible collective directions on that
domain, construct its complement or constraint Jacobian, and compare the
interacting and one-body time evolutions with the same state normalization.
The phase of the Gaussian square root must then be continued from an
explicit reference. These are genuine outstanding field-theory requirements
\Ld. They cannot be supplied by choosing $I$ or $J$ because one satisfies
an algebraic identity.

\section{Checks, corrected derivatives and failures}
For either direction, the signed momentum derivatives are
\begin{equation}
 P'_{\rm good}(a)=\frac23\epsilon\sin a\sin3a,\qquad
 P'_{\rm bad}(b)=\frac23\epsilon\sinh b\sinh3b.
\end{equation}
Both are checked symbolically, including the factor $\epsilon$ omitted
from formulas presented for both directions in the earlier report.

\begin{longtable}{p{0.51\linewidth}p{0.07\linewidth}p{0.32\linewidth}}
\toprule Result & Label & Check\\\midrule
Tail-independent collective Wronskian determinant and Darboux density
 & (c) & Exact symbolic matrices\\
Collision-time realization of that matrix
 & (c) & Four full spatial quadratures, 55 digits\\
Fourth-order symplectic boundary flux
 & (c) & Exact differential identity\\
Basis-independent collective projector
 & (c) & Exact canonical embedding\\
Fixed-endpoint obstruction and fixed-$M$ rank loss
 & (c) & Explicit localized test and symbolic restriction\\
Boundary-condition and phase controls
 & (c) & Exact finite matrices and oscillator eigenstate\\
Projected field determinant and absolute phase
 & (d) & Not computed; domain and propagation data remain\\
\bottomrule
\end{longtable}

\paragraph{Failures recorded.}
A PDF rebuild initially failed because removing a manual equation tag left
a blank paragraph inside a display; removing that blank line fixed the build.
The initial auxiliary six-dimensional canonical embedding had the opposite
sign, producing $Z^TJZ=-W$ instead of $W$. The exact matrix check caught this;
the embedding signs were corrected. No tolerance was changed. No collision
quadrature failed. The fixed-endpoint failure of the projector is an intended
audit finding, not a numerical failure. The earlier sign omission in the
momentum derivatives is corrected explicitly above.

\paragraph{References and provenance.}
This note uses the repository's tau solution and prior classical symplectic
form as stated conventions and independently checks the needed derivatives,
Wronskians and matrix identities. It adds no external literature claims or
new references. The preceding determinant audit records the extent of the
FK source passages actually read.
\end{document}
